% GENERATED FILE. Edit canonical lessons, then run npm run generate:books.
% canonical-source-hash: fnv1a64:f0fb260c5b06cf58
% canonical-lessons: FR-C267-je-pourrais, FR-C267-ce-serait, FR-C267-revision-conditionnel, FR-W267-six-lignes

\chapter{Would — the conditionnel, and Six Lines}
\label{ch:fr-would}
\hlchaptermodality{267}

\begin{chapteropening}
\textbf{By the end of this chapter:} \emph{I can say what would happen with the conditionnel, and write six connected lines about a journey that end looking forward.}

Complete FR-W267-six-lignes: six lines to one reader: the scene in the imparfait, events in order, a reason, and a last line with j'espère, ce serait or si.
\end{chapteropening}

\section[je pourrais, j'aimerais, je devrais]{je pourrais, j'aimerais — the would of French}

\label{lesson:FR-C267-je-pourrais}



\begin{quote}
Ask a stranger for a coffee. (\emph{Je voudrais un café.}) Then say \textquotedblleft{}I should sleep\textquotedblright{}. (\emph{Je devrais dormir.})

Two forms of one tense. This lesson gives it its name and its recipe.
\end{quote}

\begin{grammarlens}[title={Grammar Lens: the infinitive, plus the imparfait's endings}]
This is the \textbf{conditionnel}, French's \textquotedblleft{}would\textquotedblright{}. Keep the \textbf{whole} infinitive and add the endings of the imparfait, \textbf{-ais} for \textquotedblleft{}I\textquotedblright{}:

\begin{quote}
parler, \textbf{je parlerais} — I would speak finir, \textbf{je finirais} — I would finish prendre, \textbf{je prendrais} — I would take
\end{quote}

A verb in \emph{-re} lets its final \emph{e} go first: \emph{prendr-}, \emph{prendrais}. The other persons change only the ending: \textbf{je parlerais}, \textbf{tu parlerais}, \textbf{il parlerait}, \textbf{nous parlerions}, \textbf{vous parleriez}, \textbf{ils parleraient}. Beside the imparfait: \emph{je parlais}, I was speaking; \emph{je parlerais}, I would speak. The infinitive inside it is the whole difference.

A few common verbs squeeze their stem first. You have met two of them, \emph{vouloir} in \textbf{je voudrais} and \emph{devoir} in \textbf{je devrais}; the third is \emph{pouvoir}, \textbf{je pourrais}. \emph{Aimer} follows the plain recipe: \textbf{j'aimerais}.

Three jobs, all softening. A \textbf{polite request}: \emph{Vous pourriez m'aider ?} — Could you help me? A \textbf{wish}: \emph{J'aimerais partir en vacances.} — I'd like to go on holiday. And \textbf{advice}: \emph{À ta place, je dormirais.} — In your place, I'd sleep.
\end{grammarlens}

\begin{cousinweb}
Latin had no would-tense. Spoken Latin made one by gluing a past of \emph{habēre}, \textquotedblleft{}to have\textquotedblright{}, onto the infinitive: roughly \textquotedblleft{}I had to speak\textquotedblright{}, which slid into \textquotedblleft{}I would speak\textquotedblright{}. French chose the \textbf{imperfect} of that verb, the same \emph{habēbat} that became \emph{avait}; worn down on the end of a verb, it is the \textbf{-ait} of \emph{parlerait}. Italian glued on a different past, which is why its endings do not match. Neither EtymDB nor Lewis \& Short covers a change this late; this is the usual explanation in histories of Romance, given as that. The name \textbf{conditionnel} comes from Medieval Latin \emph{conditiōnālis}, built on \emph{conditiō}, \textquotedblleft{}a condition\textquotedblright{} (EtymDB 2.1).
\end{cousinweb}

\subsection*{Your turn}
Say these aloud:

\begin{itemize}
\raggedright
  \item \emph{je parlerais, tu parlerais, il parlerait}
  \item the imparfait, then the conditionnel — \emph{je parlais}, then \emph{je parlerais}
  \item a polite request, after a problem — \emph{Excusez-moi, j'ai un problème. Vous pourriez m'aider ?}
  \item a wish — \emph{J'aimerais partir en vacances.}
  \item an answer that depends — \emph{Ça dépend. J'aimerais venir, mais je ne peux pas.}
\end{itemize}

\begin{tcolorbox}[breakable,colback=teal!4,colframe=teal!35!black,title={Before you move on}]
Build \textquotedblleft{}I would finish\textquotedblright{}. (\emph{\textbf{Je finirais}} — the whole infinitive, plus \emph{-ais}.) What are the three jobs of the conditionnel? (A \textbf{polite request}, a \textbf{wish}, \textbf{advice}.) Which Latin past hides in its endings? (\emph{\textbf{Habēbat}}, \textquotedblleft{}he had\textquotedblright{}, the same one inside \emph{\textbf{avait}}.)
\end{tcolorbox}
\section[ce serait]{ce serait — it would be}

\label{lesson:FR-C267-ce-serait}



\begin{quote}
Build \textquotedblleft{}I would speak\textquotedblright{}. (\emph{Je parlerais.}) Then \textquotedblleft{}it's good\textquotedblright{}, now. (\emph{C'est bien.})
\end{quote}

\subsection*{You'll want to know: ce serait}
\textbf{ce serait} — \textquotedblleft{}it would be\textquotedblright{}, \textquotedblleft{}that would be\textquotedblright{}.

\textbf{être} refused the imparfait's recipe and gave you \emph{c'était}. Its conditionnel has a stem of its own too, \textbf{ser-}, with the same endings: \emph{je serais}, \emph{tu serais}, \emph{il serait}. \emph{C'est} is now, \emph{c'était} was, and \emph{\textbf{ce serait}} is the judgement on something that has not happened:

\begin{quote}
\emph{\textbf{Ce serait} bien.} — That would be good. \emph{\textbf{Ce serait} plus rapide en train.} — It would be quicker by train.
\end{quote}

\emph{En train}, \textquotedblleft{}by train\textquotedblright{}, and \emph{en voiture}, \textquotedblleft{}by car\textquotedblright{}, are how French says the way you travel.

\begin{grammarlens}[title={Grammar Lens: si, the imparfait, and the would}]
Now the condition that is \textbf{not} true. Put the \emph{si} half in the \textbf{imparfait}, and the result in the \textbf{conditionnel}:

\begin{quote}
\emph{\textbf{Si j'avais} le temps, \textbf{je partirais} en vacances.} — If I had the time, I'd go on holiday. \emph{\textbf{Si c'était} moins cher, \textbf{je prendrais} le train.} — If it were cheaper, I'd take the train.
\end{quote}

Set it beside last chapter's real condition. \emph{S'il pleut, je reste}: it might rain, and here is the plan. \emph{S'il pleuvait, je resterais}: it is not raining, and I am imagining. The conditionnel goes in the result, and the \emph{si} half takes the imparfait.

In relaxed speech, above all in Québec, you will hear \emph{si j'aurais} for \textquotedblleft{}if I had\textquotedblright{}; standard French keeps the imparfait after \emph{si}, and so should a letter or an exam.
\end{grammarlens}

\begin{cousinweb}
\textbf{être} is inherited, through Old French \emph{estre}, from Latin \emph{esse}, \textquotedblleft{}to be\textquotedblright{} (EtymDB 2.1). The stem \textbf{ser-} is usually traced to a spoken Latin infinitive \emph{essere}, reshaped from \emph{esse}, with the worn \emph{habēbat} of the last lesson on its end: something like \textquotedblleft{}had to be\textquotedblright{}. Neither source records forms this late, so that step is the usual explanation rather than a cited one.
\end{cousinweb}

\subsection*{Your turn}
Say these aloud:

\begin{itemize}
\raggedright
  \item \emph{je serais, tu serais, il serait}
  \item \emph{Ce serait plus rapide en train.}
  \item \emph{Si j'avais le temps, je partirais en vacances.}
  \item real, then imagined — \emph{S'il pleut, je reste.} then \emph{S'il pleuvait, je resterais.}
  \item a judgement with \emph{mieux} — \emph{Ce serait mieux en train.}
\end{itemize}

\begin{tcolorbox}[breakable,colback=teal!4,colframe=teal!35!black,title={Before you move on}]
What stem does \emph{être} use in the conditionnel? (\emph{\textbf{Ser-.}}) Say \textquotedblleft{}it would be quicker by train\textquotedblright{}. (\emph{\textbf{Ce serait plus rapide en train.}}) In \emph{si j'avais le temps, je partirais}, which half takes the conditionnel? (The \textbf{result}, \emph{je partirais}.)
\end{tcolorbox}
\section[(review)]{Review — if this, then that would}

\label{lesson:FR-C267-revision-conditionnel}



\begin{quote}
Say \textquotedblleft{}that would be good\textquotedblright{}. (\emph{Ce serait bien.}) Then \textquotedblleft{}I could\textquotedblright{}. (\emph{Je pourrais.})
\end{quote}

\begin{grammarlens}[title={What you've built}]
Two chapters, and a whole way of talking about what has not happened yet. A real branch: \emph{s'il pleut, je reste}. A branch not yet named: \emph{ça dépend}, or \emph{ça dépend de toi}. The softened result, in the conditionnel: \emph{je pourrais}, \emph{j'aimerais}, \emph{je devrais}, \emph{ce serait}. And the condition that is not true now: \emph{si j'avais le temps, je partirais}.

\begin{quote}
— \emph{On va au parc demain ?} — \emph{Ça dépend. Ce serait bien, mais s'il pleut, on pourrait aller au musée.}
\end{quote}

\emph{On pourrait} is \textquotedblleft{}we could\textquotedblright{}, the \emph{on} form of \emph{je pourrais}: a soft suggestion, which is why it fits after \emph{s'il pleut}. The \emph{si} half stays in the present, because rain tomorrow is a real possibility.
\end{grammarlens}

\subsection*{Your turn}
Say these aloud:

\begin{itemize}
\raggedright
  \item \emph{Ce serait bien d'aller au parc demain.}
  \item \emph{S'il pleut, on pourrait aller au musée.}
  \item advice inside a branch — \emph{Si tu es fatigué, tu devrais dormir.}
  \item a hope and a judgement — \emph{J'espère partir en vacances. Ce serait bien.}
  \item the imagined one — \emph{Si j'avais le temps, je partirais en vacances.}
\end{itemize}

\begin{tcolorbox}[breakable,colback=teal!4,colframe=teal!35!black,title={Before you move on}]
In \emph{s'il pleut, on pourrait aller au musée}, which half is the conditionnel? (\emph{\textbf{On pourrait aller.}}) Which tense does the \emph{si} half take when the condition is not true now? (The \textbf{imparfait}.) What would \textquotedblleft{}we would be\textquotedblright{} be, from \emph{serait}? (\emph{\textbf{Nous serions.}}) And which Latin verb hides in the \emph{devoir} of \emph{je devrais}? (\emph{\textbf{Dēbēre}}, to owe.)
\end{tcolorbox}
\section[(six lignes)]{Six lines about a journey}

\label{lesson:FR-W267-six-lignes}



\begin{quote}
Chapter forty-eight asked for four lines that held together, with \emph{parce que}. This one asks for six, and they tell a story.

Pick one person who will read it: \emph{un ami}, \emph{une amie}, someone in your family.
\end{quote}

\subsection*{What to know first}
A told story is not a list of events. It has a \textbf{scene}, in the imparfait: \emph{c'était}, \emph{il y avait}, \emph{il faisait froid}, \emph{j'étais content}. It has \textbf{events}, in the passé composé, put in order with \emph{d'abord}, \emph{ensuite} and \emph{enfin}. It has at least one \textbf{reason}, pointing back with \emph{parce que} or forward with \emph{donc}, and perhaps a turn with \emph{mais}.

And a good one ends by looking forward: \emph{j'espère …}, or \emph{ce serait …}, or a condition with \emph{si}.

The DELF B1 writing paper asks for a connected text of at least 160 words. Six lines is the first step towards it, not the whole of it.

\subsection*{Writing — connected composition}
On paper, write \textbf{six lines} to your reader about a journey, real or invented, one line for each of these:

\begin{enumerate}
\raggedright
  \item the scene: when it was, and the weather;
  \item where you were, who was with you, and how long you stayed, with \emph{pendant};
  \item the first thing that happened, with \emph{d'abord};
  \item the next thing, with \emph{ensuite}, and a reason, with \emph{parce que} or \emph{donc};
  \item how it ended and how you felt, with \emph{enfin} and \emph{j'étais};
  \item one line forward, with \emph{j'espère}, \emph{ce serait} or \emph{si}.
\end{enumerate}

Use only words you already write. The exercise is choosing the right past for each line.

\begin{tcolorbox}[breakable,colback=teal!4,colframe=teal!35!black,title={Before you move on}]
Now go through your lines as you went through the journey in chapter two hundred and fifty-nine. Mark every verb \textbf{S} for scene or \textbf{E} for event. Lines one and two should be mostly \textbf{S}; three and four mostly \textbf{E}; five has one of each. A line about the weather marked \textbf{E} is the commonest slip: change \emph{il a plu} back to \emph{il pleuvait}.

Then two spelling checks. The participle of an \emph{-er} verb ends in \textbf{-é}: \emph{j'ai parlé}, \emph{je suis arrivé}, and with \emph{être} it agrees with a woman writing about herself: \emph{je suis arrivée}. And the imparfait's endings are silent but spelled: \emph{-ais}, \emph{-ait}, \emph{-aient}.

Last, read line six on its own. If it has a \emph{si} about something real, is the verb after it in the plain present? If it has \emph{ce serait}, is it your judgement on something that has not happened yet? Then it looks forward, and the story is finished.
\end{tcolorbox}
